\section{Introduction}

Symbolic music generation represents music as structured symbolic events
\citep{huangPopMusicTransformer2020,fradetMiditokPythonPackage2023}, such
as note onsets, durations, velocities, programs, and timing information, rather
than as raw audio waveforms. This representation makes it possible to model
musical structure with sequence models while keeping the output editable as MIDI.
Instead of predicting samples in the audio domain, the system predicts symbolic
events that can later be decoded into notes, instruments, tempo changes, and
other performance-related information.

In this project, MIDI files are converted into discrete token sequences and
modeled with autoregressive neural language models. This follows the same broad
sequence-modeling view used by recent symbolic-music Transformer systems
\citep{hsiaoCompoundWordTransformer2021,ensMMMExploringConditional2020}. The
central question of the thesis is how an xLSTM model compares with
Transformer-based language models when all three are trained for the same
symbolic music generation task. Rather than treating the models as isolated
implementations, the work uses a shared experimental pipeline so that xLSTM,
GPT-2-style, and LLaMA-style models can be compared under matched data,
tokenization, context-length, sampling, and evaluation conditions.

The repository\footnote{\url{https://github.com/Alkhatir/duet-of-models}} is
organized around this controlled comparison. The first stage develops the data
and modelling infrastructure: selecting a usable MIDI subset, constructing
leakage-aware train/validation/test splits, choosing a REMI+-style token
representation, training xLSTM models, tracking validation behaviour, and
selecting a strong xLSTM configuration. The second stage reuses this derived
setup for Transformer experiments, training GPT-2-style and LLaMA-style models
on the same representation and evaluating their generated MIDI outputs against
the xLSTM baseline.

The thesis also includes a practical interface for using the trained models. A
Gradio application is developed to load selected checkpoints, accept or choose
MIDI prompts, generate continuations with configurable sampling parameters, and
present the resulting MIDI files for inspection and playback. This is important
because symbolic music generation cannot be judged from scalar metrics alone:
validation loss and token accuracy describe modelling behaviour, while the
generated MIDI files reveal musical continuity, instrumentation, rhythmic
coherence, and failure cases such as sparse or fragmented generations.

The main goal of this thesis is therefore to document and evaluate a complete
symbolic-music generation workflow rather than a single model in isolation. The
following sections describe the dataset and split construction, preprocessing
and tokenization choices, model architectures, training setup, hyperparameter
search, model selection, comparative results, limitations, and the interactive
generation interface. Together, these components support a reproducible
comparison between recurrent and Transformer-based sequence models for symbolic
music generation.
